\documentclass[11pt,letterpaper]{article}

\usepackage[top=1in, bottom=1in, left=1in, right=1in, headheight=14pt]{geometry}
\usepackage{fontspec}
\setmainfont{DejaVu Serif}
\setsansfont{DejaVu Sans}
\setmonofont{DejaVu Sans Mono}

\usepackage{xcolor}
\usepackage{titlesec}
\usepackage{fancyhdr}
\usepackage{enumitem}
\usepackage{hyperref}
\usepackage{booktabs}
\usepackage{tabularx}
\usepackage{longtable}
\usepackage{colortbl}
\usepackage{multirow}
\usepackage{array}
\usepackage{parskip}
\usepackage{tcolorbox}
\usepackage{graphicx}
\usepackage{lastpage}

% ── Color Scheme: History / Music industry (burgundy + gold + slate) ─────────
\definecolor{primary}{HTML}{4A148C}
\definecolor{secondary}{HTML}{F57F17}
\definecolor{tertiary}{HTML}{37474F}
\definecolor{accent}{HTML}{6A1B9A}
\definecolor{lightprimary}{HTML}{F3E5F5}
\definecolor{lightsecondary}{HTML}{FFF8E1}
\definecolor{lighttertiary}{HTML}{ECEFF1}
\definecolor{rowgray}{HTML}{F5F5F5}
\definecolor{bordergray}{HTML}{BDBDBD}
\definecolor{subtlegray}{HTML}{757575}
\definecolor{darktext}{HTML}{212121}

% ── Section Formatting ────────────────────────────────────────────────────────
\titleformat{\section}
  {\Large\bfseries\sffamily\color{primary}}
  {\thesection}{1em}{}
  [\vspace{-0.5em}\textcolor{bordergray}{\rule{\textwidth}{0.4pt}}]

\titleformat{\subsection}
  {\large\bfseries\sffamily\color{secondary}}
  {\thesubsection}{1em}{}

\titleformat{\subsubsection}
  {\normalsize\bfseries\sffamily\color{tertiary}}
  {\thesubsubsection}{1em}{}

\titlespacing*{\section}{0pt}{1.5em}{0.8em}
\titlespacing*{\subsection}{0pt}{1.2em}{0.5em}
\titlespacing*{\subsubsection}{0pt}{0.8em}{0.3em}

% ── Custom Environments ───────────────────────────────────────────────────────
\newcommand{\stageheader}[2]{%
  \noindent\colorbox{#2}{\makebox[\textwidth][l]{%
    \textcolor{white}{\bfseries\sffamily\hspace{6pt}#1\hspace{6pt}}}}%
  \vspace{0.5em}\par%
}

\newtcolorbox{asciibox}{
  colback=rowgray, colframe=bordergray,
  arc=1pt, boxrule=0.5pt,
  left=6pt, right=6pt, top=4pt, bottom=4pt,
  fontupper=\small\ttfamily,
  breakable
}

\newtcolorbox{quotebox}{
  colback=lightprimary, colframe=primary,
  arc=2pt, boxrule=0.5pt,
  left=12pt, right=12pt, top=8pt, bottom=8pt,
  breakable
}

% ── Table helpers ─────────────────────────────────────────────────────────────
\newcolumntype{L}[1]{>{\raggedright\arraybackslash}p{#1}}
\newcolumntype{C}[1]{>{\centering\arraybackslash}p{#1}}
\newcommand{\tableheader}[1]{\textbf{\sffamily\textcolor{white}{#1}}}
\newcommand{\metafield}[2]{\textbf{\sffamily #1} #2\par}

% ── Headers / Footers ────────────────────────────────────────────────────────
\pagestyle{fancy}
\fancyhf{}
\fancyhead[L]{\small\sffamily\textcolor{subtlegray}{Module 05 Split Analysis}}
\fancyhead[R]{\small\sffamily\textcolor{subtlegray}{GSD Mission Package}}
\fancyfoot[C]{\small\sffamily\textcolor{subtlegray}{Page \thepage\ of \pageref{LastPage}}}
\renewcommand{\headrulewidth}{0.4pt}
\renewcommand{\footrulewidth}{0pt}

% ── Hyperlinks ────────────────────────────────────────────────────────────────
\hypersetup{
  colorlinks=true,
  linkcolor=secondary,
  urlcolor=secondary,
  pdfauthor={GSD Ecosystem},
  pdftitle={Module 05 Split Analysis -- Mission Package},
  pdfsubject={Vision to Mission Pipeline},
}

% ═════════════════════════════════════════════════════════════════════════════
\begin{document}

% ── Title Page ───────────────────────────────────────────────────────────────
\thispagestyle{empty}
\begin{center}
\vspace*{3cm}

{\small\sffamily\textcolor{primary}{\textbf{GSD RESEARCH MISSION PACKAGE}}}

\vspace{0.3cm}
\textcolor{bordergray}{\rule{8cm}{0.4pt}}

\vspace{1.5cm}
{\fontsize{26}{32}\selectfont\bfseries\sffamily\textcolor{primary}{Module 05 Split Analysis\\[0.3em]NastyMix Records \& CD Place Memory}}

\vspace{0.8cm}
{\Large\sffamily\textcolor{secondary}{Seattle Hip-Hop Cultural Series}}

\vspace{0.5cm}
{\large\sffamily\itshape\textcolor{tertiary}{What Could Be Better --- A Structural Remediation Study}}

\vspace{3cm}
{\sffamily\textcolor{accent}{Vision $\rightarrow$ Research $\rightarrow$ Mission}}\\[0.3em]
{\small\sffamily\textcolor{accent}{Full Pipeline Execution}}

\vspace{3cm}
{\small\sffamily\textcolor{subtlegray}{Date: March 2026  |  Status: Ready for GSD Orchestrator}}\\[0.3em]
{\small\sffamily\textcolor{subtlegray}{Activation Profile: Squadron  |  Estimated: 8--10 context windows across 2--3 sessions}}
\end{center}
\newpage

% ── TOC ──────────────────────────────────────────────────────────────────────
\tableofcontents
\newpage

% ═════════════════════════════════════════════════════════════════════════════
\stageheader{STAGE 1 --- VISION DOCUMENT}{primary}

\section{Module 05 Split Analysis --- Vision Guide}

\metafield{Date:}{March 2026}
\metafield{Status:}{Structural Remediation / Cold-Start Research}
\metafield{Parent Series:}{Seattle Hip-Hop Cultural Study}
\metafield{Problem:}{Module 05 at 613 lines is the longest single module in the series}
\metafield{Depends on:}{Central District Cultural Study (CD mission), deepaudioanalyzermission.pdf}
\metafield{Context:}{Investigating whether Module 05 should be decomposed into two focused modules}

\subsection{Vision}

A record label is not the same thing as the neighborhood that produced it. NastyMix Records was
born in the Central District, but it quickly became something else: a business, a brand, a distribution
machine, a legal battleground. By 1990, Nastymix had signed acts from Tacoma, New Jersey, Miami,
and a British Euro-disco duo --- not because Ed Locke stopped caring about Seattle, but because the
label's logic had outgrown any single neighborhood. The music industry story and the place story began
diverging the moment ``Square Dance Rap'' broke in Flint, Michigan.

Module 05's 613 lines are overloaded because they are trying to tell both stories at once. The NastyMix
roster --- Sir Mix-A-Lot, Kid Sensation, High Performance, Criminal Nation, and a dozen lesser-known
signings --- belongs to a music industry narrative about independent label economics, regional distribution,
the cassette economy, and what happens when a West Coast label's anchor artist walks out the door. The
Central District mapping, MOHAI documentation, and gentrification analysis belong to a place narrative
about a neighborhood that went from 65\% Black to under 20\% Black in a single generation, and the
cultural institutions that are racing to document what remains.

These are not the same story. They share geography and a time period, but their analytical frames,
primary sources, and questions of consequence diverge sharply. The label story asks: how does an
independent hip-hop label work, what does it build, and what destroys it? The place story asks: what
happens to a neighborhood after the music moves out?

Separating them produces two modules of approximately 280--320 lines each --- consistent with the
series median and far more focused in execution. It also reveals a structural insight: the NastyMix
label story is a music industry module, while the CD place memory module is a cultural geography and
archival module. They require different research methodologies, different primary sources, and different
analytical frames. Keeping them fused was always a category error.

\subsection{Problem Statement}

\begin{enumerate}
\item \textbf{Category mismatch in a single module.} NastyMix Records is a music industry subject
(label economics, roster strategy, distribution deals, artist-label disputes). The Central District's
gentrification arc is a cultural geography subject (spatial mapping, demographic data, archival
documentation, community memory). These require different analytical frames, different primary
sources, and different success criteria. Fusing them into Module 05 produced a 613-line overshoot
that serves neither subject adequately.

\item \textbf{The label roster is underdocumented at the individual artist level.} The NastyMix
roster extended well beyond Sir Mix-A-Lot: Kid Sensation, High Performance (Tacoma), Criminal
Nation (Tacoma), The Acu\"sed (Oak Harbor metal), Adrienne (R\&B), Whiz Kid, Bob \& the Mob
(New Jersey), Side F-X (Miami), Rodney O \& Joe Cooley (California). Many of these artists are
poorly documented in existing modules because they were compressed into a roster list rather than
receiving the contextual treatment their role in the label's economics demands.

\item \textbf{The MOHAI documentation layer is architecturally important and currently thin.}
MOHAI's ``Legacy of Seattle Hip-Hop'' exhibit (September 2015 -- May 2016) is the single most
important institutional documentation of the CD hip-hop scene, curated by community members
Jazmyn Scott and Aaron Walker-Loud and recognized by the AASLH with a 2016 Leadership in
History Award. This exhibit warrants detailed treatment as a primary archival source, not a
passing citation.

\item \textbf{The gentrification data deserves its own analytical arc.} The Central District went
from over 65\% Black in the 1970s--80s to under 20\% Black by 2020 (MOHAI/AASLH documentation).
This is a well-documented displacement event with its own timeline, causal factors (redlining,
tech boom, rising rents), artistic responses (Draze's ``The Hood Ain't The Same,'' 2014), and
institutional preservation efforts (NAAM, 206 Zulu, Washington Hall, Wa Na Wari). Treating it
as a sub-point within a label history module buries its significance.

\item \textbf{The split creates two well-scoped parallelizable modules.} Module 05a (NastyMix
Label Ecosystem) and Module 05b (Central District Place Memory) can run simultaneously in the
series wave plan, each with a clear self-contained scope, distinct primary sources, and specific
success criteria. The split also aligns Module 05a with the music industry cluster and Module
05b with the community heritage cluster within the larger series architecture.
\end{enumerate}

\subsection{Core Concept}

\textbf{One Zip Code, Two Stories: Separate What Diverges, Connect What Converges}

The NastyMix label and the Central District neighborhood share a founding moment --- the Rotary
Boys and Girls Club, Sir Mix-A-Lot, Nes Rodriguez, Ed Locke, a Commodore 64, and a cassette.
But from that shared origin, they follow divergent trajectories. The label expands outward (Tower
Building, Gold Party at the Four Seasons, national distribution, a New Jersey signing, a British
duo). The neighborhood contracts inward under gentrification pressure. Separating the two modules
follows the logic of the divergence. Connecting them through the series wave plan preserves the
shared origin.

\subsubsection{Module Architecture After Split}

\begin{asciibox}
BEFORE SPLIT
  Module 05 (613 lines)
  |-- NastyMix roster (label economics)
  |-- Central District mapping (geography)
  |-- MOHAI documentation (archives)
  `-- Gentrification analysis (displacement)
  [CATEGORY ERROR: music industry + cultural geography fused]

AFTER SPLIT
  Module 05a: NastyMix Label Ecosystem        (~300 lines)
  |-- Founding story (1985, four co-founders)
  |-- Complete roster with contextual profiles
  |-- Business arc: $900 startup to platinum to collapse
  |-- Distribution mechanics and cassette economics
  |-- Radio pipeline: KFOX FreshTracks as label infrastructure
  `-- Collapse: Mix departs 1991, rights, Ichiban partnership

  Module 05b: Central District Place Memory   (~300 lines)
  |-- Spatial mapping: neighborhood boundaries and key sites
  |-- MOHAI as primary archive (Legacy exhibit 2015-16)
  |-- Gentrification arc: 65%+ Black -> under 20% (1970-2020)
  |-- Redlining foundation: restrictive covenants, segregation
  |-- Artistic documentation: Draze "The Hood Ain't The Same"
  `-- Preservation infrastructure: NAAM, 206 Zulu, Wa Na Wari

SERIES INTEGRATION
  05a feeds: Music industry cluster (label economics, artist development)
  05b feeds: Community heritage cluster (MOHAI, NAAM, spatial memory)
  Both connect to: Module 03 (CD ecosystem) as shared foundation
\end{asciibox}

\subsection{Research Modules}

\subsubsection{MA: NastyMix Founding and Co-Founder Profiles}

The label's founding is a four-person story often collapsed into two. Ed Locke (former KYAC DJ,
$900 budget, near bankruptcy on first pressing) and Greg Jones (investor) are the business side.
Nes Rodriguez (Filipino-American, FreshTracks host, the ``Nasty'' of the name) and Sir Mix-A-Lot
(Anthony Ray, the artist, the product) are the music side. The Tower Building offices at 1809 7th
Avenue Suite 800 represent the physical expression of the label's success. The NAMA Outstanding
Achievement Award and back-to-back ``Best Company of the Year'' and ``Independent Record Label
of the Year'' awards (1990) mark its institutional recognition. Locke's stated goal of expanding
from rap to ``a variety of pop styles'' explains the diverse and ultimately fatal late-period signings.

\subsubsection{MB: Full Roster --- Regional Acts}

The Northwest core: Sir Mix-A-Lot (Auburn/CD), Kid Sensation (Stephen Spence, Mix's teen prot\'eg\'e,
first appearance on \textit{Swass} 1988, debut LP \textit{Rollin' With Number One} 1990), High Performance
(Tacoma hard-core rap, both dance tunes and political songs), Criminal Nation (Tacoma, active 1990--92,
\textit{Release The Pressure} 1990, \textit{Trouble In The Hood} 1992, DJ Eugenius as regional connector).
The Acu\"sed (Oak Harbor speed-metal, the label's ill-fated rock expansion). These Northwest acts share
geographic proximity and the KFOX/FreshTracks promotional pipeline.

\subsubsection{MC: Full Roster --- National Signings}

Locke's 1990 expansion: Adrienne Baseden (R\&B vocalist, Wilberforce University scholar, Whitney
Houston-vein), Whiz Kid (former New York rapper), Bob \& the Mob (New Jersey, MC Speedo and
DJ Slayer, funk-rap updates of Spinners and Marvin Gaye), Side F-X (Miami, Reggie ``Kut Mighty
Swift'' La Lanne), Rodney O \& Joe Cooley (California). These signings represent Locke's attempt
to diversify beyond Seattle rap --- and the strategic error that diluted the label's identity and
ultimately failed to produce any second-tier successes.

\subsubsection{MD: Business Arc and Collapse}

\$900 startup budget (1985) $\rightarrow$ ``Square Dance Rap'' breaks in Flint, Michigan $\rightarrow$
SWASS platinum (1988) $\rightarrow$ Gold Party at the Four Seasons Olympic Hotel (April 29, 1989)
$\rightarrow$ Seminar gold (1989) $\rightarrow$ Mix departs 1991, takes all rights $\rightarrow$
Ichiban partnership (rocky) $\rightarrow$ closure 1992. The structural vulnerability: a single-artist
label whose anchor artist owned his masters. When Mix left, he took the revenue base. The Northwest
rap expansion (Kid Sensation as ``Seattle's next rap star'') failed to produce a second platinum act.

\subsubsection{ME: CD Spatial Mapping and Key Sites}

The Central District's hip-hop geography: Bryant Manor apartments (19th Ave and East Yesler Way,
Mix's childhood residence), Rotary Boys and Girls Club (CD, founding meeting point), Rainier Vista
Boys and Girls Club (South Seattle, Mix's early DJ parties), Martin Luther King Jr. Way (formerly
Empire Way, as noted in Mix's lyrics), 23rd Ave and Union (Uncle Ike's controversy, Draze's
``Irony on 23rd Street''), MLK Jr. Way and Cherry Street (The Facts newspaper building, converted
to dog daycare --- gentrification's emblem in Draze's documentation).

\subsubsection{MF: MOHAI as Primary Archive}

``The Legacy of Seattle Hip-Hop'' (MOHAI, September 2015 -- May 2016): curated by Jazmyn Scott
and Aaron Walker-Loud; partnership with the Black Heritage Society; the February 2014 community
day drew nearly 1,000 attendees (``wholly unexpected''); eight-month exhibit covering graffiti,
DJing, dance, production, emceeing; Dr.\ Daudi Abe lecture series; 2016 AASLH Leadership in
History Award; 2016 Charles Redd Center for Western Studies Honorable Mention. MOHAI's
``Segregated Seattle: From Redlining to Gentrification'' (programming, in partnership with NAAM).
MOHAI digital photograph collection (civil rights documentation). MOHAI as Smithsonian Affiliate.

\subsubsection{MG: Gentrification Arc and Preservation Response}

Displacement timeline: African Americans arrive Seattle 1880s--90s (coal mine labor recruitment),
settle East Madison and Yesler--Jackson $\rightarrow$ Central District consolidates as Black cultural
hub $\rightarrow$ 1970s--80s: over 65\% Black (AASLH/MOHAI data) $\rightarrow$ 2015: CD designated
Arts \& Cultural District $\rightarrow$ 2020: under 20\% Black. Causal factors: redlining (restrictive
covenants documented by UW Civil Rights \& Labor History Project), tech-economy rent pressure,
school busing (1978--1999, Mix-A-Lot documented this directly). Artistic response: Draze's ``The
Hood Ain't The Same'' (2014), ``Irony on 23rd Street.'' Preservation infrastructure: Northwest
African American Museum (NAAM), 206 Zulu, Washington Hall, Wa Na Wari (``Our Home,'' co-founded
2019 by Inye Wokoma et al., former Black-owned family home as art space).

\subsection{Chipset Configuration}

\begin{asciibox}
# chipset.yaml -- Module 05 Split
chipset:
  name: "mod05-split-analysis"
  version: "1.0.0"
  description: "Module 05 structural decomposition: NastyMix label + CD place memory"
  parent_series: "seattle-hip-hop-cultural-study"

modules:
  mod05a:
    name: "nastymix-label-ecosystem"
    scope: "Music industry: founding, roster, economics, collapse"
    primary_source: "HistoryLink.org (files 9793, 9794), Daudi Abe Emerald Street"
    line_target: 280-320
    tracks: [MA, MB, MC, MD]
    cluster: "music-industry"

  mod05b:
    name: "central-district-place-memory"
    scope: "Cultural geography: spatial mapping, MOHAI archives, gentrification arc"
    primary_source: "MOHAI Legacy exhibit, UW Civil Rights & Labor History Project"
    line_target: 280-320
    tracks: [ME, MF, MG]
    cluster: "community-heritage"

integration:
  shared_foundation: "Module 03 (Central District ecosystem)"
  cross_reference: "05a <-> 05b at founding moment (Rotary Boys & Girls Club)"
  series_wave: "Both 05a and 05b run in Wave 1 as parallel tracks"
\end{asciibox}

\subsection{Scope Boundaries}

\subsubsection{In Scope (Module 05a --- NastyMix Label Ecosystem)}

\begin{itemize}
\item Complete four-founder story (Locke, Jones, Rodriguez, Ray) with individual profiles
\item Full roster with contextual treatment of each act's geographic origin and role
\item Business arc from \$900 budget to platinum to collapse, including the Ichiban partnership
\item The Gold Party (1989) and Fifth Anniversary Party (1990) as cultural milestones
\item Artist-label dispute mechanics (Mix's departure, rights ownership)
\item KFOX FreshTracks as the label's de facto promotional infrastructure
\item NAMA awards as institutional recognition data points
\item The failed rock/pop expansion as a structural analysis case
\end{itemize}

\subsubsection{In Scope (Module 05b --- Central District Place Memory)}

\begin{itemize}
\item CD spatial mapping with specific addresses and neighborhood boundaries
\item MOHAI ``Legacy of Seattle Hip-Hop'' exhibit as primary archival documentation
\item Gentrification arc with demographic data (65\%+ to under 20\% Black, 1970--2020)
\item Redlining and restrictive covenant history (UW Civil Rights \& Labor History Project)
\item School busing program (1978--1999) as spatial and social context
\item Draze's ``The Hood Ain't The Same'' (2014) as artistic documentation source
\item Preservation infrastructure: NAAM, 206 Zulu, Washington Hall, Wa Na Wari
\item 2015 Arts \& Cultural District designation
\end{itemize}

\subsubsection{Out of Scope (Both Modules)}

\begin{itemize}
\item Detailed musical analysis (covered by audio analysis modules in series)
\item Post-Nastymix Sir Mix-A-Lot career (Rhyme Cartel, Def American, Mack Daddy)
\item New NastyMix Records revival (Maurice ``MC Duce'' Owens, 2020s)
\item National gentrification policy analysis (CD-specific only)
\end{itemize}

\subsection{Success Criteria}

\begin{enumerate}
\item Module 05a lands within 280--320 lines, consistent with series median.
\item Module 05b lands within 280--320 lines, consistent with series median.
\item Every artist on the NastyMix roster receives a minimum two-sentence contextual entry
   (origin, role, and fate post-label), not merely a name in a list.
\item The business arc in 05a is traceable as a timeline: founding year $\rightarrow$ first
   national hit $\rightarrow$ Gold Party $\rightarrow$ Fifth Anniversary $\rightarrow$ Mix
   departure $\rightarrow$ Ichiban $\rightarrow$ closure, with dates and sales figures sourced
   to HistoryLink.org and Emerald Street.
\item MOHAI's Legacy exhibit in 05b is documented at the level of: curators, duration, community
   partnerships, attendee count (the 2014 community day and the 8-month exhibit), and AASLH
   award year.
\item The gentrification arc in 05b includes at least three demographic data points with sources,
   covers the redlining foundation, and names specific institutions in the preservation response.
\item Both modules explicitly cross-reference at the Rotary Boys and Girls Club founding moment
   (the shared origin point before the label story and place story diverge).
\item The split is validated by the series wave plan: 05a and 05b can execute as parallel tracks
   in Wave 1 with no inter-module dependency.
\item All roster claims in 05a cross-reference against Daudi Abe's \textit{Emerald Street}
   (UW Press, 2020) as the primary scholarly authority.
\item The Nasty Nes preservation urgency marker (health crisis documented Nov. 2024, Seattle Times)
   is retained in 05a as a provenance note for first-person source accessibility.
\end{enumerate}

\subsection{Relationship to Other Documents}

\begin{center}
\rowcolors{2}{rowgray}{white}
\begin{tabularx}{\textwidth}{L{4cm}XL{3cm}}
\toprule
\rowcolor{primary}
\tableheader{Document} & \tableheader{Relationship} & \tableheader{Direction} \\
\midrule
Central District Cultural Study & Foundation module for the CD ecosystem; 05b deepens the place memory layer it established & Extends \\
deepaudioanalyzermission.pdf & Parent series context; Module 05 overshoot identified in What Could Be Better review & Parent \\
gsd-foxfire-heritage-skills-vision.md & Parallel community documentation methodology (oral history, cultural geography) & Lateral \\
Emerald Street (Daudi Abe, UW Press 2020) & Primary scholarly authority for all roster and scene claims in both modules & Source \\
HistoryLink.org (files 9793, 9794) & Primary source for NastyMix milestone events and party documentation & Source \\
MOHAI Legacy of Seattle Hip-Hop (2015) & Primary archival source for 05b; community-curated exhibit with AASLH award & Source \\
UW Civil Rights \& Labor History Project & Primary source for redlining, restrictive covenants, and gentrification data & Source \\
\bottomrule
\end{tabularx}
\end{center}

\subsection{The Through-Line}

The NastyMix label and the Central District share a 1985 origin point, diverge almost immediately,
and are connected in 2015 at MOHAI by Jazmyn Scott and Aaron Walker-Loud --- two community curators
who understood that the label's story and the neighborhood's story are inseparable even when they
require separate analytical frames.

That is the Amiga Principle applied to historical research: you can run the label economics and the
cultural geography in parallel, each with its own specialized execution path, because the architecture
of the larger series holds them in relation. They converge in the origin (the Rotary Boys and Girls
Club) and they converge in the preservation moment (the Legacy exhibit). In between, they do their
own work, ask their own questions, and produce their own findings. That is not a bug in Module 05's
original design. It is the module telling you it needs two chips.

% ═════════════════════════════════════════════════════════════════════════════
\newpage
\stageheader{STAGE 2 --- RESEARCH REFERENCE}{secondary}

\section{Module 05 Split Analysis --- Research Reference}

\subsection{How to Use This Document}

This section provides sourced findings for each of the seven research tracks (MA--MG) organized
across the two proposed modules. Tracks MA--MD feed Module 05a (NastyMix Label Ecosystem).
Tracks ME--MG feed Module 05b (Central District Place Memory). The primary scholarly authority
for all scene claims is Daudi Abe, \textit{Emerald Street: A History of Hip-Hop in Seattle}
(University of Washington Press, 2020).

Key source organizations:
\begin{itemize}
\item \textbf{HistoryLink.org} --- Seattle/Washington state online encyclopedia; primary source for
   NastyMix milestone documentation (files 9793, 9794)
\item \textbf{MOHAI (Museum of History \& Industry)} --- Primary archival repository for Central
   District hip-hop documentation; Smithsonian Affiliate, AASLH award-winner
\item \textbf{UW Civil Rights \& Labor History Project} --- Redlining, restrictive covenants,
   gentrification bibliography and primary document archive
\item \textbf{Humanities Washington / Spark Magazine} --- Daudi Abe's contextualization of the
   NastyMix era for general audiences (January 2022)
\item \textbf{Seattle Times Archive} --- Period reporting on Nastymix Records (June 1990,
   October 1992) and Nasty Nes health crisis (November 2024)
\item \textbf{AASLH (American Association of State and Local History)} --- MOHAI Legacy exhibit
   documentation, Leadership in History Award 2016
\end{itemize}

\subsection{Track MA: NastyMix Founding Profiles}

\subsubsection{The Four Co-Founders}

\begin{center}
\rowcolors{2}{rowgray}{white}
\begin{tabularx}{\textwidth}{L{3.5cm}L{3cm}X}
\toprule
\rowcolor{primary}
\tableheader{Co-Founder} & \tableheader{Background} & \tableheader{Role at Nastymix} \\
\midrule
Ed Locke & Former KYAC DJ; self-described ``disco freak''; acted in high school / UW drama & CEO; business operations; relationship management with radio, press, distributors \\
Greg Jones & Investor & Initial capital (label launched on \$900 budget); near-bankruptcy on first pressing \\
``Nasty'' Nes Rodriguez & Filipino-American; arrived US 1970; KFOX FreshTracks host from 1980 & Promotions; radio pipeline; the ``Nasty'' of NastyMix; left Seattle 1997 for LA (HITS magazine) \\
Anthony ``Sir Mix-A-Lot'' Ray & Born Auburn 1963; grew up CD; Bryant Manor apartments & Sole recording artist at founding; creative engine; left 1991 with all masters \\
\bottomrule
\end{tabularx}
\end{center}

The label name fused Rodriguez's ``Nasty'' with Ray's ``Mix-A-Lot.'' The label's operational address
moved from ``humble offices off Western Avenue'' to the Tower Building at 1809 7th Avenue Suite 800
--- a spatial indicator of growth. Locke acknowledged the label almost died before launch: California
pressing company oversupply, returns, near-bankruptcy, saved by Locke cold-calling radio stations
until ``Square Dance Rap'' broke in Flint, Michigan.

\textit{Source: Seattle Times archive, June 3, 1990 (``A Mix Bag''); Humanities Washington/Spark (Jan. 2022)}

\subsection{Track MB: Northwest Roster Profiles}

\subsubsection{Sir Mix-A-Lot}

Two platinum Nastymix albums: \textit{SWASS} (1988, 1M+ shipped) and \textit{Seminar} (1989, 500K+).
``Posse on Broadway'' (1988) was the record that ``put Seattle on the rap map'' (Seattle Times). Mix
wrote, arranged, programmed, performed, produced, and engineered his recordings himself. Left 1991
after dispute over creative control and promotional strategy; took all rights. Estimated \$100M+ lifetime
royalties from ``Baby Got Back'' alone. \textit{Source: HistoryLink file 9793; Seattle Times; Blackpast.org}

\subsubsection{Kid Sensation (Stephen Spence)}

Mix's teen prot\'eg\'e; first appearance on \textit{Swass} (1988). Solo debut \textit{Rollin' With Number One}
(Nastymix, 1990). Second LP \textit{The Power of Rhyme} (Nastymix/Ichiban/Emerald City Records, 1992,
featuring Ken Griffey Jr. --- ``The Way We Swing''). Left Nastymix for Ichiban; released \textit{Seatown Funk}
(1995). Later produced for Jackie Jackson; created jingles for NFL and NBA franchises. Considered a
Northwest legend. \textit{Source: Good Ol'Dayz; Town Love/Crane City Music}

\subsubsection{High Performance}

Tacoma hard-core rap group. Group leader: Maurice ``MC Duce'' Owens (later revived the Nastymix
brand in 2020s). Characterized by both dance tunes and ``hard political songs'' (Seattle Times, 1990).
Performed at the 1989 Gold Party and 1990 Fifth Anniversary Party. Represented the label's South
Sound geographic reach. \textit{Source: HistoryLink file 9794; New Nastymix Records}

\subsubsection{Criminal Nation}

Tacoma-based; active 1990--1992 on Nastymix. Albums: \textit{Release The Pressure} (1990),
\textit{Trouble In The Hood} (1992). Briefly resurfaced 1998 (Ocean Records, \textit{Resurrection}).
DJ Eugenius was a regional connector who later worked with Sir Mix-A-Lot and other Northwest artists
as writer, DJ, and producer. Wojack (MC Deff) released two post-Nastymix solo albums (1997, 2001).
\textit{Source: Good Ol'Dayz}

\subsubsection{The Acu\"sed}

Oak Harbor (Whidbey Island) speed-metal band. Nastymix's most conspicuous genre departure.
Announced at the 1990 Fifth Anniversary Party alongside the rap acts. Represented Locke's attempt
to diversify the label beyond Black music and beyond Seattle, which The Rocket and industry observers
questioned. \textit{Source: HistoryLink file 9794; Seattle Times, June 3, 1990}

\subsection{Track MC: National Signings}

By 1990, Locke doubled the label's client list with acts from outside the Northwest:

\begin{center}
\rowcolors{2}{rowgray}{white}
\begin{tabularx}{\textwidth}{L{3cm}L{2.5cm}X}
\toprule
\rowcolor{primary}
\tableheader{Act} & \tableheader{Origin} & \tableheader{Nastymix Role / Notes} \\
\midrule
Adrienne (Baseden) & Wilberforce, OH & R\&B vocalist, Whitney Houston style; full LP delayed by college schedule \\
Whiz Kid & New York & Street-oriented rapper; Nastymix debut singles only \\
Bob \& the Mob & New Jersey & MC Speedo and DJ Slayer; funk-rap covers of Spinners and Marvin Gaye \\
Side F-X & Miami & Reggie ``Kut Mighty Swift'' La Lanne; Nastymix release in late 1990 \\
Rodney O \& Joe Cooley & California & Listed on final Nastymix roster; never prominent in Seattle scene \\
Rococo & United Kingdom & Euro-disco twin sisters; Locke's furthest genre reach \\
\bottomrule
\end{tabularx}
\end{center}

None of the national signings produced charting releases. The strategic rationale was diversification;
the practical result was dilution of the label's Northwest identity and financial exposure without
offsetting revenue. \textit{Source: Seattle Times, June 3, 1990}

\subsection{Track MD: Business Arc and Collapse}

\subsubsection{Milestone Timeline}

\begin{center}
\rowcolors{2}{rowgray}{white}
\begin{tabularx}{\textwidth}{L{1.5cm}L{3cm}X}
\toprule
\rowcolor{primary}
\tableheader{Year} & \tableheader{Milestone} & \tableheader{Significance} \\
\midrule
1985 & Label founded & \$900 budget; Ed Locke + Greg Jones + Nes Rodriguez + Sir Mix-A-Lot \\
1985 & ``Square Dance Rap'' EP & 45,000+ copies sold; breaks in Flint, MI and other scattered markets \\
1988 & \textit{SWASS} released & Platinum (1M shipped); spent 1+ year on both Billboard Black and Pop LP charts \\
Apr. 1989 & Gold Party at Four Seasons & Spanish Ballroom, Olympic Hotel; 400+ attendees; SWASS Gold award \\
1989 & \textit{Seminar} released & Gold (500K+); ``My Hooptie'' hit; Mix unhappy with promotion \\
Nov. 1990 & Fifth Anniversary Party & Phantom Club (332 5th Ave N); SWASS Platinum + Seminar Gold presented; NAMA awards \\
1990--91 & National expansion & Roster doubles; rock, pop, R\&B, international signings; none charted \\
1991 & Mix departs & Took all masters and rights; Rhyme Cartel formed with Ricardo Frazer \\
1991--92 & Ichiban partnership & Georgia-based distributor; described by Locke as ``rocky'' \\
Late 1992 & Label closes & Kid Sensation's 1992 LP carried Ichiban distribution logo; final release \\
\bottomrule
\end{tabularx}
\end{center}

\textit{Source: HistoryLink files 9793, 9794; Seattle Times October 6, 1992; Town Love/Crane City Music}

\subsubsection{Structural Vulnerability}

The label's fatal structural problem: its entire revenue base resided in a single artist who owned
his masters. When Mix departed, he took everything --- catalog rights, future royalties, and the
promotional credibility that had underwritten the national expansion. Kid Sensation was positioned
as ``Seattle's next rap star'' but never approached Mix's commercial performance. The national
signings produced no revenue to compensate.

\subsection{Track ME: Central District Spatial Mapping}

\subsubsection{Key Sites}

\begin{center}
\rowcolors{2}{rowgray}{white}
\begin{tabularx}{\textwidth}{L{4.5cm}X}
\toprule
\rowcolor{primary}
\tableheader{Location} & \tableheader{Significance} \\
\midrule
Bryant Manor Apartments, 19th Ave \& E. Yesler Way & Sir Mix-A-Lot's childhood residence; recording origin point (Commodore 64 setup) \\
Rotary Boys and Girls Club, Central District & Founding meeting point: Mix meets Nes Rodriguez; NastyMix origin story \\
Rainier Vista Boys and Girls Club, South Seattle & Mix's early DJ party venue (pre-Nastymix, 1983) \\
Martin Luther King Jr. Way (formerly Empire Way) & Central axis of the Black CD; Mix references original name in ``What's Real'' \\
23rd Ave \& Union & Uncle Ike's pot shop controversy (2010s gentrification flashpoint; Draze ``Irony on 23rd St'') \\
MLK Jr. Way \& Cherry Street & Site of \textit{The Facts} newspaper (Seattle's oldest Black paper, founded 1961); converted to dog daycare \\
Washington Hall & Historic CD cultural venue; preservation anchor \\
\bottomrule
\end{tabularx}
\end{center}

\textit{Source: Sir Mix-A-Lot Wikipedia; Emerald Street excerpts (artenoir.org); Seattle Globalist}

\subsection{Track MF: MOHAI as Primary Archive}

\subsubsection{The Legacy of Seattle Hip-Hop Exhibit}

\begin{itemize}
\item \textbf{Dates:} September 19, 2015 -- May 1, 2016 (8 months)
\item \textbf{Curators:} Jazmyn Scott (owner, The Town Entertainment; co-founder 50 Next: Seattle
   Hip-Hop Worldwide) and Aaron Walker-Loud (founder, Big World Breaks; co-founder 50 Next;
   Education Director, Seattle JazzED)
\item \textbf{Partnership:} MOHAI + Black Heritage Society
\item \textbf{Origin:} February 2014 community day; nearly 1,000 attendees (``wholly unexpected'')
   triggered the full 8-month exhibit
\item \textbf{Content:} Audio recordings, photography, artwork, artifacts; graffiti, DJing, dance,
   production, emceeing; first-person narratives; QTPOC community panel (DJ Sassy Black Cat, moderator)
\item \textbf{Dr.\ Daudi Abe lecture:} Previewed themes of \textit{Emerald Street}; covered the
   Emerald Street Boys, FreshTracks, NastyMix era
\item \textbf{Awards:} 2016 AASLH Leadership in History Award; 2016 Charles Redd Center for
   Western Studies Exhibition Excellence Honorable Mention
\end{itemize}

\textit{Source: MOHAI exhibit page; AASLH.org case study; Seattle Times September 18, 2015}

\subsubsection{MOHAI's Broader CD Documentation Role}

``Segregated Seattle: From Redlining to Gentrification'' (MOHAI programming, partnership with
NAAM; sold-out event). MOHAI digital photograph collection (civil rights documentation, indexed in
UW Civil Rights \& Labor History Project bibliography). MOHAI as Smithsonian Affiliate. The AASLH
case study notes the exhibit ``became a kind of community center'' for a neighborhood undergoing
rapid gentrification.

\subsection{Track MG: Gentrification Arc and Preservation Response}

\subsubsection{Demographic Timeline}

\begin{center}
\rowcolors{2}{rowgray}{white}
\begin{tabularx}{\textwidth}{L{2cm}X}
\toprule
\rowcolor{primary}
\tableheader{Period} & \tableheader{Status} \\
\midrule
1880s--90s & African Americans arrive Seattle for coal mine labor; settle East Madison and Yesler--Jackson \\
Early 20th c. & Restrictive covenants confine Black residents to the Central Area \\
1960s--80s & Over 65\% of Central District residents are Black (MOHAI/AASLH documentation) \\
1978--1999 & School busing program: minorities bused from CD and South End to North End schools \\
2014 & Draze releases ``The Hood Ain't The Same'' video; CD gentrification documented in hip-hop \\
2015 & CD designated an official Arts \& Cultural District \\
2020 & Under 20\% of Central District residents are Black (MOHAI/AASLH documentation) \\
\bottomrule
\end{tabularx}
\end{center}

\subsubsection{Artistic Documentation: Draze}

Dumi ``Draze'' Maraire Jr. (b. Seattle; parents: Dumisani Abraham Maraire Sr., Zimbabwean musician;
childhood time in Zimbabwe). ``The Hood Ain't The Same'' (2014): shot around the CD and South End;
documents the conversion of \textit{The Facts} newspaper building to a dog daycare; ``we was red lined
in but now we black balled out so they can sell green.'' ``Irony on 23rd Street'': Uncle Ike's pot shop
at 23rd \& Union (``how many brothers went to jail on this corner from moving dime bags...'').
Video screened at MOHAI as part of the Legacy exhibit programming.

\textit{Source: Arte Noir (Emerald Street excerpts); Seattle Globalist; TheGrio}

\subsubsection{Preservation Infrastructure}

\begin{center}
\rowcolors{2}{rowgray}{white}
\begin{tabularx}{\textwidth}{L{3.5cm}L{1.8cm}X}
\toprule
\rowcolor{primary}
\tableheader{Institution} & \tableheader{Founded} & \tableheader{Role} \\
\midrule
NAAM (Northwest African American Museum) & 2008 & Primary museum for Black Seattle cultural heritage; MOHAI partner \\
206 Zulu & Est. & Seattle hip-hop cultural organization; referenced in MOHAI exhibit as preservation anchor \\
Washington Hall & Historic & Historic CD performance venue; preservation anchor for Black performance arts \\
Wa Na Wari (``Our Home'') & 2019 & Former Black-owned family home (Frank and Goldyne Green's); co-founded by Inye Wokoma et al.; art exhibits, performances, screenings \\
\bottomrule
\end{tabularx}
\end{center}

\textit{Source: The Root (7 Black Sites); AASLH Legacy exhibit documentation}

\subsection{Safety and Sensitivity Considerations}

\begin{center}
\rowcolors{2}{rowgray}{white}
\begin{tabularx}{\textwidth}{L{3.5cm}L{2cm}X}
\toprule
\rowcolor{primary}
\tableheader{Area} & \tableheader{Type} & \tableheader{Handling} \\
\midrule
Artist-label dispute details & GATE & Present Mix-A-Lot's departure as documented fact; do not speculate on financial terms beyond what is sourced \\
Gentrification causality & ANNOTATE & Present evidence (redlining, busing, tech economy) without policy advocacy \\
Nes Rodriguez health status & ANNOTATE & Flag as preservation urgency (first-person source accessibility at risk); use respectfully, not as tragedy narrative \\
Cultural ownership claims & GATE & ``Our Home'' (Wa Na Wari), Black Heritage Society --- community-asserted framing; honor without editorializing \\
QTPOC community content & GATE & From Legacy exhibit panels; present without over-summarizing or flattening community-led framing \\
\bottomrule
\end{tabularx}
\end{center}

\subsection{Source Bibliography}

\subsubsection{Primary Archival Sources}
\begin{itemize}
\item Daudi Abe. \textit{Emerald Street: A History of Hip-Hop in Seattle 1979--2015.} University of Washington Press, 2020. [Primary scholarly authority]
\item HistoryLink.org, File 9793: ``Nastymix Records party celebrates Gold Record awarded to Sir Mix-A-Lot's SWASS album, April 29, 1989.''
\item HistoryLink.org, File 9794: ``Nastymix Records hosts fifth Anniversary party, November 29, 1990.''
\item Seattle Times Archive, June 3, 1990: ``A Mix Bag --- Seattle's Nastymix Records Is Becoming a Key Player...'' (period reporting)
\item Seattle Times Archive, October 6, 1992: ``Nastymix Records To Close.''
\item Seattle Times, November 27, 2024: ``Community rallies around Seattle rap producer, DJ who shaped local scene.'' [Nes Rodriguez health crisis]
\end{itemize}

\subsubsection{Institutional Documentation}
\begin{itemize}
\item MOHAI. ``The Legacy of Seattle Hip-Hop'' exhibit page. mohai.org/exhibits/the-legacy-of-seattle-hip-hop/
\item AASLH. ``The Legacy of Seattle Hip-Hop'' case study. aaslh.org/the-legacy-of-seattle-hip-hop/ (2018)
\item Daudi Abe excerpts at Arte Noir. artenoir.org/post/excerpts-from-emerald-street (July 2025)
\item UW Civil Rights \& Labor History Project. ``African American Civil Rights Movement History in Seattle: Bibliography.'' depts.washington.edu/civilr/bib\_af\_am.htm
\end{itemize}

\subsubsection{Secondary Sources}
\begin{itemize}
\item Humanities Washington / Spark. ``How the Northwest Shocked the Hip-hop World.'' January 20, 2022.
\item Blackpast.org. ``Anthony `Sir Mix-A-Lot' Ray (1963--).''
\item AASLH case study. ``The Legacy of Seattle Hip-Hop.'' (2018)
\item Seattle Globalist. ``What the South End Sounds Like: Draze Drops `Seattle's Own'.'' March 31, 2016.
\item TheGrio. ``Hip-hop artist mourns gentrification of black Seattle neighborhood.'' February 28, 2018.
\item The Root. ``7 Black Sites You Must See in Seattle.''
\item New Nastymix Records. ``About Us.'' newnastymixrecords.com/about-us/
\item Town Love / Crane City Music. Sir Mix-A-Lot, Kid Sensation, Criminal Nation archives. townlove.cranecitymusic.com
\item Good Ol'Dayz. Kid Sensation, Criminal Nation artist pages. thegoodoldayz.com
\end{itemize}

% ═════════════════════════════════════════════════════════════════════════════
\newpage
\stageheader{STAGE 3 --- MISSION PACKAGE}{tertiary}

\section{Milestone Specification}

\metafield{Mission Objective:}{Decompose Module 05 of the Seattle Hip-Hop Cultural Series into two focused modules:
Module 05a (NastyMix Label Ecosystem, music industry frame) and Module 05b (Central District Place Memory,
cultural geography frame). Each module targets 280--320 lines, consistent with series median.
The split resolves the category-error overshoot and enables parallel execution.}

\subsection{Architecture Overview}

\begin{asciibox}
WAVE 0: Foundation
  [Haiku] Shared schemas, cross-reference contracts, series integration map

WAVE 1: Parallel Execution
  [Sonnet] Track A: 05a Label Ecosystem (MA+MB+MC+MD)
  [Sonnet] Track B: 05b Place Memory (ME+MF+MG)
  [Both run simultaneously -- no inter-module dependency until Wave 2]

WAVE 2: Synthesis
  [Opus] Cross-module integration: shared founding moment, series arc validation
         Emerald Street cross-reference audit (all roster claims)
         Line count calibration (target: 280-320 each)

WAVE 3: Verification + Publication
  [Sonnet] Final modules delivered to series; success criteria checklist; CAPCOM review
\end{asciibox}

\subsection{System Layers}

\begin{enumerate}
\item \textbf{Source Layer} --- HistoryLink, Emerald Street, MOHAI exhibit, UW CRL\&HP, Seattle Times archive
\item \textbf{Research Layer} --- Seven tracks (MA--MG) producing structured findings per module
\item \textbf{Module Layer} --- 05a (music industry) and 05b (cultural geography) as distinct deliverables
\item \textbf{Series Integration Layer} --- Cross-reference contracts linking both modules to Module 03 and series wave plan
\end{enumerate}

\subsection{Deliverables}

\begin{center}
\rowcolors{2}{rowgray}{white}
\begin{tabularx}{\textwidth}{C{0.5cm}L{4cm}XL{3cm}}
\toprule
\rowcolor{primary}
\tableheader{\#} & \tableheader{Deliverable} & \tableheader{Acceptance Criteria} & \tableheader{Spec Track} \\
\midrule
1 & Module 05a: NastyMix Label Ecosystem & 280--320 lines; full roster; business arc with dates; Emerald Street cross-ref & MA, MB, MC, MD \\
2 & Module 05b: Central District Place Memory & 280--320 lines; MOHAI exhibit documented; gentrification arc with 3+ data points & ME, MF, MG \\
3 & Cross-reference contract & Both modules cite shared founding moment; 05a $\leftrightarrow$ 05b link documented & Wave 2 \\
4 & Series integration update & Series wave plan updated: 05a and 05b as parallel Wave 1 tracks replacing original 05 & Wave 3 \\
\bottomrule
\end{tabularx}
\end{center}

\subsection{Component Breakdown}

\begin{center}
\rowcolors{2}{rowgray}{white}
\begin{tabularx}{\textwidth}{L{4cm}L{2.5cm}C{1.5cm}C{2cm}}
\toprule
\rowcolor{primary}
\tableheader{Component} & \tableheader{Dependencies} & \tableheader{Model} & \tableheader{Est. Tokens} \\
\midrule
Wave 0: Shared schemas & None & Haiku & 5K \\
Track MA: Founding profiles & Wave 0 & Sonnet & 12K \\
Track MB: Northwest roster & Wave 0, MA & Sonnet & 15K \\
Track MC: National signings & Wave 0, MA & Sonnet & 10K \\
Track MD: Business arc & Wave 0, MA--MC & Sonnet & 12K \\
Track ME: Spatial mapping & Wave 0 & Sonnet & 12K \\
Track MF: MOHAI archive & Wave 0, ME & Sonnet & 15K \\
Track MG: Gentrification arc & Wave 0, ME, MF & Sonnet & 15K \\
Wave 2: Synthesis + audit & MA--MG & Opus & 25K \\
Wave 3: Verification & All & Sonnet & 10K \\
\bottomrule
\end{tabularx}
\end{center}

\subsection{Wave Execution Plan}

\subsubsection{Wave Summary}

\begin{center}
\rowcolors{2}{rowgray}{white}
\begin{tabularx}{\textwidth}{L{1.2cm}XC{2cm}C{2cm}C{2cm}}
\toprule
\rowcolor{primary}
\tableheader{Wave} & \tableheader{Work} & \tableheader{Parallel} & \tableheader{Model} & \tableheader{Est. Tokens} \\
\midrule
0 & Foundation schemas & No & Haiku & 5K \\
1 & 05a tracks (MA--MD) + 05b tracks (ME--MG) & Yes (2 tracks) & Sonnet & 64K \\
2 & Synthesis, Emerald Street audit, calibration & No & Opus & 25K \\
3 & Verification, series integration update & No & Sonnet & 10K \\
\midrule
\multicolumn{4}{r}{\textbf{Total estimated:}} & \textbf{$\sim$104K} \\
\bottomrule
\end{tabularx}
\end{center}

\subsubsection{Dependency Graph}

\begin{asciibox}
Wave 0 ──> Shared Schemas
            |
            +──> Track A (05a): MA -> MB -> MC -> MD ──+
            |                                           |
            +──> Track B (05b): ME -> MF -> MG ────────+
                                                        |
                                               Wave 2: Synthesis (Opus)
                                                        |
                                               Wave 3: Verification (Sonnet)
                                                        |
                                               Series Integration Update
\end{asciibox}

\subsection{Test Plan}

\subsubsection{Test Categories}

\begin{center}
\rowcolors{2}{rowgray}{white}
\begin{tabularx}{\textwidth}{L{3.5cm}C{1.5cm}L{2cm}X}
\toprule
\rowcolor{primary}
\tableheader{Category} & \tableheader{Count} & \tableheader{Priority} & \tableheader{Action on Failure} \\
\midrule
Safety-Critical & 4 & Mandatory & BLOCK: do not deliver \\
Scholarly Authority & 4 & Mandatory & BLOCK: do not deliver \\
Core Functionality & 8 & High & Revise before delivery \\
Integration & 4 & Medium & Note and flag \\
\bottomrule
\end{tabularx}
\end{center}

\subsubsection{Safety-Critical Tests (BLOCK)}

\begin{center}
\rowcolors{2}{rowgray}{white}
\begin{tabularx}{\textwidth}{L{1.5cm}L{4cm}X}
\toprule
\rowcolor{primary}
\tableheader{ID} & \tableheader{Verifies} & \tableheader{Expected Result} \\
\midrule
SC-01 & No undocumented financial claims & All dollar figures and sales numbers attributed to HistoryLink or Seattle Times or Abe \\
SC-02 & Nes Rodriguez health note handled appropriately & Note included as preservation urgency, not exploited as tragedy narrative \\
SC-03 & Gentrification data sourced & All demographic percentages attributed to MOHAI/AASLH documentation \\
SC-04 & QTPOC Legacy exhibit content & Panel content presented from community-curated framing; not flattened or summarized beyond MOHAI documentation \\
\bottomrule
\end{tabularx}
\end{center}

\subsubsection{Scholarly Authority Tests (BLOCK)}

\begin{center}
\rowcolors{2}{rowgray}{white}
\begin{tabularx}{\textwidth}{L{1.5cm}L{4cm}X}
\toprule
\rowcolor{primary}
\tableheader{ID} & \tableheader{Verifies} & \tableheader{Expected Result} \\
\midrule
SA-01 & Emerald Street cross-reference & All formation-era roster claims cross-referenced against Abe (2020) before finalization \\
SA-02 & HistoryLink primary sourcing & Gold Party (1989) and Fifth Anniversary (1990) details sourced to files 9793 and 9794 specifically \\
SA-03 & MOHAI exhibit accuracy & Exhibit dates, curators, partnerships, and awards match MOHAI exhibit page and AASLH documentation \\
SA-04 & UW Civil Rights project for redlining & Redlining and restrictive covenant claims reference UW CRL\&HP bibliography \\
\bottomrule
\end{tabularx}
\end{center}

\subsubsection{Verification Matrix}

\begin{center}
\rowcolors{2}{rowgray}{white}
\begin{tabularx}{\textwidth}{L{5.5cm}L{2.5cm}C{2cm}}
\toprule
\rowcolor{primary}
\tableheader{Success Criterion} & \tableheader{Test ID(s)} & \tableheader{Status} \\
\midrule
Both modules land 280--320 lines & CF-01, CF-02 & \\
Full roster with 2+ sentence profiles & CF-03 & \\
Business arc traceable as dated timeline & CF-04, SA-02 & \\
MOHAI exhibit documented at full detail & CF-05, SA-03 & \\
Gentrification arc with 3+ sourced data points & CF-06, SC-03 & \\
Founding moment cross-reference in both modules & INT-01 & \\
Wave plan updated with 05a/05b as parallel tracks & INT-02 & \\
All roster claims cross-referenced to Emerald Street & SA-01 & \\
Nes Rodriguez preservation note handled appropriately & SC-02 & \\
Both modules pass CAPCOM review before commit & CF-07 & \\
\bottomrule
\end{tabularx}
\end{center}

\subsection{Mission Crew Manifest}

\begin{center}
\rowcolors{2}{rowgray}{white}
\begin{tabularx}{\textwidth}{L{2.5cm}L{3cm}X}
\toprule
\rowcolor{primary}
\tableheader{Role} & \tableheader{Agent} & \tableheader{Responsibility} \\
\midrule
Mission Commander & DIRECTOR & Overall split execution; wave sequencing; series integration \\
Label Historian & ARCHIVIST-A & Tracks MA--MD; NastyMix founding, roster, business arc \\
Place Historian & ARCHIVIST-B & Tracks ME--MG; spatial mapping, MOHAI, gentrification \\
Synthesis Lead & WEAVER & Wave 2 cross-module integration; founding moment cross-reference \\
Source Auditor & VERITAS & Emerald Street cross-reference audit; HistoryLink sourcing verification \\
Safety Warden & GUARDIAN & SC-01 through SC-04; community content sensitivity review \\
Series Integrator & NAVIGATOR & Series wave plan update; 05a/05b integration into larger series \\
CAPCOM & CAPCOM & Human checkpoint before final commit \\
\bottomrule
\end{tabularx}
\end{center}

\subsection{Execution Summary}

\begin{center}
\rowcolors{2}{rowgray}{white}
\begin{tabularx}{\textwidth}{L{5cm}X}
\toprule
\rowcolor{primary}
\tableheader{Metric} & \tableheader{Value} \\
\midrule
Total research tracks & 7 (MA--MG) \\
Parallel execution tracks & 2 (Wave 1: 05a + 05b simultaneously) \\
Sequential depth & 4 waves \\
Model split & Opus $\sim$25\% / Sonnet $\sim$65\% / Haiku $\sim$10\% \\
Estimated context windows & $\sim$8--10 \\
Estimated sessions & 2--3 \\
Total estimated tokens & $\sim$104K \\
Total tests & 20 (8 mandatory-pass BLOCK) \\
Target line count per module & 280--320 (series median) \\
Primary scholarly authority & Daudi Abe, \textit{Emerald Street} (UW Press, 2020) \\
\bottomrule
\end{tabularx}
\end{center}

\vspace{2em}
\begin{quotebox}
\textit{``Nastymix's 1989 `Gold Party' was not only the moment which defined Northwest Hip-Hop's
arrival as a force in the local music community \ldots''}

\vspace{0.5em}
{\small\textcolor{subtlegray}{--- Glen Boyd, cited in HistoryLink file 9793}}

\vspace{1em}
The label and the neighborhood share a founding myth. Separating them into two modules
does not diminish that myth --- it gives each half room to breathe, and room to mean
what it actually means. The NastyMix story is a story about what independent labels
can build from a \$900 budget and a cassette player. The Central District story is a
story about what a community preserves after the buildings change hands.
Both stories need the full stage.
\end{quotebox}

\vspace{1em}
\begin{center}
{\small\sffamily\textcolor{subtlegray}{GSD Research Mission Package $\cdot$ Module 05 Split Analysis $\cdot$ March 2026 $\cdot$ Ready for GSD Orchestrator}}
\end{center}

\end{document}
